\subsection{\texorpdfstring{Example: linear representations in the spaces \(L^{p}\)}{Example: linear representations in the spaces L\^{}\{p\}}}

Again let G be a locally compact group, operating continuously on the left in a locally compact space X. Let \(\beta\) be a positive measure on X with support X. Let us assume that there exists a continuous function \(\chi > 0\) on \(G \times X\) such that, for every \(s \in G\) ,

\[
\boldsymbol {\gamma} (s) \beta = \chi (s ^ {- 1}, \cdot) \cdot \beta
\]

(which implies in particular that \(\beta\) is quasi-invariant under \(\mathbf{G}\) ). Then, \(\chi\) is a multiplier. For, given \(s, t\) in \(\mathbf{G}\) , one has

\[
\begin{array}{r l} \boldsymbol {\gamma} (s) \boldsymbol {\gamma} (t) \beta = \boldsymbol {\gamma} (s) \big (\chi (t ^ {- 1}, \cdot) \cdot \beta \big) & = \big (\boldsymbol {\gamma} (s) \chi (t ^ {- 1}, \cdot) \big) \cdot \big (\boldsymbol {\gamma} (s) \beta \big) \\ & = \big (\boldsymbol {\gamma} (s) \chi (t ^ {- 1}, \cdot) \big) \cdot \chi (s ^ {- 1}, \cdot) \cdot \beta , \\ \boldsymbol {\gamma} (s t) \beta & = \chi (t ^ {- 1} s ^ {- 1}, \cdot) \cdot \beta , \end{array}
\]

therefore

\[
\chi (t ^ {- 1}, s ^ {- 1} x) \chi (s ^ {- 1}, x) = \chi (t ^ {- 1} s ^ {- 1}, x)
\]

locally \(\beta\) -almost everywhere, consequently everywhere, since \(\chi\) is continuous and \(\beta\) has support \(X\) .

Let \(p \in [1, +\infty[\) . For every \(f \in \mathcal{L}_{\mathbf{C}}^{p}(\mathbf{X}, \beta)\) and every \(s \in \mathbf{G}\) , let \(\gamma_{\chi,p}(s)f\) be the function on \(\mathbf{X}\) defined by

\[
\left(\gamma_ {\chi , p} (s) f\right) (x) = \chi (s ^ {- 1}, x) ^ {1 / p} f (s ^ {- 1} x).
\]

One has

\[
\begin{array}{c} \int^ {*} | \chi (s ^ {- 1}, x) ^ {1 / p} f (s ^ {- 1} x) | ^ {p} d \beta (x) = \int^ {*} | f (s ^ {- 1} x) | ^ {p} \chi (s ^ {- 1}, x) d \beta (x) \\ = \int | f (x) | ^ {p} d \beta (x), \end{array}
\]

therefore \(\gamma_{\chi,p}(s)f\in\mathcal{L}_{\mathbf{C}}^{p}(\mathrm{X},\beta)\) . One sees that \(\gamma_{\chi,p}(s)\) is an isometric endomorphism of \(\mathcal{L}_{\mathbf{C}}^{p}(\mathrm{X},\beta)\) and defines, by passage to the quotient, an isometric endomorphism of \(\mathrm{L}_{\mathbf{C}}^{p}(\mathrm{X},\beta)\) , also denoted \(\gamma_{\chi,p}(s)\) . On the other hand, \(\chi^{1/p}\) is obviously a multiplier, therefore \(\gamma_{\chi,p}\) is a linear representation of G in \(\mathrm{L}_{\mathbf{C}}^{p}(\mathrm{X},\beta)\) by what we have seen in No.~3.

PROPOSITION 8. --- The linear representation \(\gamma_{\chi,p}\) of G in \(\mathrm{L}_{\mathbf{C}}^{p}(\mathbf{X},\beta)\) is continuous and isometric.

Let \(f \in \mathcal{K}(\mathrm{X})\) . When \(s\) tends to \(s_0\) in \(\mathrm{G}\) , \(\boldsymbol{\gamma}_{\chi,p}(s)f\) tends to \(\boldsymbol{\gamma}_{\chi,p}(s_0)f\) in \(\mathcal{K}(\mathrm{X})\) , hence in \(\mathrm{L}_{\mathbf{C}}^p (\mathrm{X},\beta)\) . Since the \(\boldsymbol{\gamma}_{\chi,p}(s)\) are isometric, Prop. 8 is obtained by applying Remark 2 of No.~1.

For the case that \(\chi\) is not assumed continuous, cf.~§4, Exer. 13.

PROPOSITION 9. --- Suppose that each function \(\chi(s,\cdot)\) is bounded. Then \(\gamma_{\chi}\) leaves \(\mathrm{L}_{\mathbf{C}}^{p}(\mathrm{X},\beta)\) stable, and the linear representation \(\gamma_{\chi}\) of \(\mathbf{G}\) in \(\mathrm{L}_{\mathbf{C}}^{p}(\mathrm{X},\beta)\) is continuous.

Let \(f\in \mathcal{L}_{\mathbf{C}}^{p}(\mathrm{X},\beta)\) . Then

\[
\begin{array}{r l} \int^ {*} | \chi (s ^ {- 1}, x) f (s ^ {- 1} x) | ^ {p} d \beta (x) \\ & \leqslant \sup _ {x \in X} \chi (s ^ {- 1}, x) ^ {p - 1} \int^ {*} | f (s ^ {- 1} x) | ^ {p} \chi (s ^ {- 1}, x) d \beta (x) \\ & = \sup _ {x \in X} \chi (s ^ {- 1}, x) ^ {p - 1} \int | f (x) | ^ {p} d \beta (x), \end{array}
\]

therefore \(\pmb{\gamma}_{\chi}(s)f\in \mathcal{L}_{\mathbf{C}}^{p}(\mathrm{X},\beta)\) , and

\[
\left\| \boldsymbol {\gamma} _ {\chi} (s) \right\| \leqslant \sup _ {x \in X} \chi (s ^ {- 1}, x) ^ {1 / q},\tag{5}
\]

where \(q\) denotes the exponent conjugate to \(p\) . If \(f \in \mathcal{K}(\mathrm{X})\) , then \(\boldsymbol{\gamma}_{\chi}(s)f\) tends to \(\boldsymbol{\gamma}_{\chi}(s_0)f\) in \(\mathcal{K}(\mathrm{X})\) , hence in \(\mathcal{L}_{\mathbf{C}}^{p}(\mathrm{X},\beta)\) , as \(s\) tends to \(s_0\) . Therefore the representation \(\boldsymbol{\gamma}_{\chi}\) of \(\mathrm{G}\) in \(\mathrm{L}_{\mathbf{C}}^{p}(\mathrm{X},\beta)\) is continuous (No.~1, Prop. 2).

Properties analogous to those of Nos. 3, 4, 5 hold if G operates on the right in X.

In particular, if one regards G as operating on itself by left or right translations, and if one takes \(\chi = 1\) , one obtains the left and right regular representations of G in \(\mathcal{C}(\mathrm{G})\) , \(\mathcal{K}(\mathrm{G})\) , \(\overline{\mathcal{K}(\mathrm{G})}\) , \(\mathcal{C}'(\mathrm{G})\) , \(\mathcal{M}(\mathrm{G})\) , \(\mathcal{M}^{1}(\mathrm{G})\) . If one takes \(\beta\) to be a left (resp. right) Haar measure on G, and if one takes \(\chi = 1\) , one obtains the left (resp. right) regular representation of G in \(\mathrm{L}_{\mathbf{C}}^{p}(\mathrm{G}, \beta)\) .

